\documentclass[aspectratio=169,11pt]{beamer}

% =============================================================================
% Template: coauthor_descriptives_template.tex
%
% Purpose:
%   Beamer template for a descriptive-results update to share with coauthors.
%   The deck is designed to compile even if some figures or tables have not yet
%   been generated: missing files are replaced with placeholders.
%
% Expected folders (relative to this .tex file):
%   figures/
%   tables/
%
% Main design choice:
%   The core deck emphasizes the full sample, reserving couple-sample tables for
%   an appendix / later discussion on within-household outcomes.
% =============================================================================

\usetheme{Madrid}
\usecolortheme{default}
\setbeamertemplate{navigation symbols}{}
\setbeamertemplate{footline}[frame number]

\usepackage{graphicx}
\usepackage{booktabs}
\usepackage{adjustbox}
\usepackage{array}
\usepackage{ragged2e}
\usepackage{appendixnumberbeamer}
\usepackage{xcolor}
\usepackage{tikz}
\usepackage{ifthen}

% ---- Paths ------------------------------------------------------------------
\newcommand{\figdir}{figures}
\newcommand{\tabdir}{tables}

% ---- Helper boxes for missing figures / tables ------------------------------
\newcommand{\missingfigurebox}[1]{%
  \begin{tikzpicture}
    \draw[rounded corners=3pt, thick, gray!70] (0,0) rectangle (10.8,5.9);
    \node[align=center, text width=9.6cm] at (5.4,2.95)
    {\small Figure not found:\\[0.3em]\texttt{#1}};
  \end{tikzpicture}%
}

\newcommand{\missingtablebox}[1]{%
  \begin{tikzpicture}
    \draw[rounded corners=3pt, thick, gray!70] (0,0) rectangle (10.8,5.1);
    \node[align=center, text width=9.6cm] at (5.4,2.55)
    {\small Table not found:\\[0.3em]\texttt{#1}};
  \end{tikzpicture}%
}

\newcommand{\safeincludegraphics}[2][]{%
  \IfFileExists{#2}{\includegraphics[#1]{#2}}{\missingfigurebox{#2}}%
}

\newcommand{\safetableinput}[1]{%
  \IfFileExists{#1}{\input{#1}}{\missingtablebox{#1}}%
}

% ---- Title info -------------------------------------------------------------
\title{COVID-19, Work Arrangements, and Household Division of Labor}
\subtitle{Descriptive update for coauthor discussion}
\author{Your Name}
\date{\today}

\begin{document}

% =============================================================================
% Title slide
% =============================================================================
\begin{frame}
  \titlepage
\end{frame}

% =============================================================================
% Roadmap
% =============================================================================
\begin{frame}{Roadmap}
\begin{itemize}
  \item Baseline grouping logic and descriptive sample overview
  \item Short-run COVID labor-market disruptions in the full sample
  \item Work outside versus work from home by baseline group
  \item Longer-run descriptive patterns in main-study follow-up waves
  \item Open questions for the empirical strategy discussion
\end{itemize}
\end{frame}

% =============================================================================
% Motivation / setup
% =============================================================================
\begin{frame}{Descriptive setup for discussion}
\begin{itemize}
  \item Current goal: summarize reduced descriptive facts before locking in the causal design.
  \item Main source of heterogeneity is the baseline 2019 industry / occupation grouping.
  \item At this stage, we mostly emphasize the full sample to see broad employment and work-arrangement trends.
  \item Couple-only objects remain useful for later discussion of household outcomes, but are not the focus of the main descriptive deck.
\end{itemize}
\vspace{0.5em}
\begin{block}{Variables we will likely treat differently in the structural / causal stage}
\begin{itemize}
  \item Baseline group / industry / occupation = pre-determined exposure.
  \item Work outside and WFH-some = realized work-arrangement margins.
  \item Time use, labor supply, earnings, health, and family outcomes = downstream outcomes.
\end{itemize}
\end{block}
\end{frame}

% =============================================================================
% Baseline sample / group table
% =============================================================================
\begin{frame}{Sample composition for later couple analysis (appendix-style reminder)}
\small
\centering
\safetableinput{\tabdir/sample_table_3x3_all.tex}
\vspace{0.5em}

\footnotesize
This is not the main focus of the descriptive deck, but it is useful to keep the eventual couple-analysis sample in view.
\end{frame}

% =============================================================================
% April 2020 snapshot
% =============================================================================
\begin{frame}{April 2020 snapshot: work status and furlough by baseline group}
\begin{columns}[T,onlytextwidth]
  \column{0.49\textwidth}
  \centering
  \safeincludegraphics[width=\linewidth,height=0.70\textheight,keepaspectratio]{\figdir/work_status_april20_groups_perc.png}

  \column{0.49\textwidth}
  \centering
  \safeincludegraphics[width=\linewidth,height=0.70\textheight,keepaspectratio]{\figdir/furlough_april20_groups_perc.png}
\end{columns}
\end{frame}

% =============================================================================
% Worked at all / hours
% =============================================================================
\begin{frame}{Short-run employment disruption: did people work at all?}
\centering
\safeincludegraphics[width=0.92\linewidth,height=0.78\textheight,keepaspectratio]{\figdir/worked_at_all_groups_overtime.png}
\end{frame}

\begin{frame}{Short-run intensive margin: hours worked by baseline group}
\centering
\safeincludegraphics[width=0.92\linewidth,height=0.78\textheight,keepaspectratio]{\figdir/hours_groups_overtime.png}
\end{frame}

% =============================================================================
% Work outside / WFH some
% =============================================================================
\begin{frame}{Work arrangement margin I: working outside the home}
\centering
\safeincludegraphics[width=0.92\linewidth,height=0.78\textheight,keepaspectratio]{\figdir/workoutside_overtime_groups.png}
\end{frame}

\begin{frame}{Work arrangement margin II: working from home at least sometimes}
\centering
\safeincludegraphics[width=0.92\linewidth,height=0.78\textheight,keepaspectratio]{\figdir/wfh_some_overtime_groups.png}
\end{frame}

\begin{frame}{Distribution of WFH intensity by baseline group}
\centering
\safeincludegraphics[width=0.96\linewidth,height=0.78\textheight,keepaspectratio]{\figdir/wfh_groups_overtime.png}
\end{frame}

% =============================================================================
% Longer-run follow-up
% =============================================================================
\begin{frame}{Longer-run follow-up: working outside by baseline group}
\centering
\safeincludegraphics[width=0.92\linewidth,height=0.78\textheight,keepaspectratio]{\figdir/future_all_workoutside_by_group_industry_based_wave_main.png}
\end{frame}

\begin{frame}{Longer-run follow-up: WFH at least sometimes by baseline group}
\centering
\safeincludegraphics[width=0.92\linewidth,height=0.78\textheight,keepaspectratio]{\figdir/future_all_wfh_some_by_group_industry_based_wave_main.png}
\end{frame}

\begin{frame}{Longer-run follow-up: weekly hours by baseline group}
\centering
\safeincludegraphics[width=0.92\linewidth,height=0.78\textheight,keepaspectratio]{\figdir/future_all_jbhrs_by_group_industry_based_wave_main.png}
\end{frame}

% =============================================================================
% Optional follow-up outcomes
% =============================================================================
\begin{frame}{Optional descriptive extension: employment status by sex}
\centering
\safeincludegraphics[width=0.96\linewidth,height=0.78\textheight,keepaspectratio]{\figdir/future_all_jbstat_facet_sex_wave_main.png}
\end{frame}

\begin{frame}{Optional descriptive extension: mental health by sex (yearly)}
\centering
\safeincludegraphics[width=0.92\linewidth,height=0.78\textheight,keepaspectratio]{\figdir/future_all_scghq1_dv_by_sex_year_main.png}
\end{frame}

% =============================================================================
% Discussion slide
% =============================================================================
\begin{frame}{Questions for coauthor discussion}
\begin{enumerate}
  \item Which descriptive margins are most informative for the eventual design: employment, work outside, WFH, hours, earnings, or family outcomes?
  \item Do the baseline group patterns look strong enough to motivate a reduced-form design built around predetermined exposure?
  \item Should the main descriptive narrative emphasize:
    \begin{itemize}
      \item sector shutdown versus key-worker contrast,
      \item work-outside versus WFH margin,
      \item or a spouse-level joint-treatment framing from the start?
    \end{itemize}
  \item Which outcomes should be in the first pass of the main paper tables once we move beyond descriptives?
\end{enumerate}
\end{frame}

% =============================================================================
% Appendix
% =============================================================================
\appendix

\begin{frame}{Appendix: detailed group version of work outside}
\centering
\safeincludegraphics[width=0.92\linewidth,height=0.78\textheight,keepaspectratio]{\figdir/workoutside_overtime_detailed_groups.png}
\end{frame}

\begin{frame}{Appendix: detailed group version of WFH-some}
\centering
\safeincludegraphics[width=0.92\linewidth,height=0.78\textheight,keepaspectratio]{\figdir/wfh_some_overtime_detailed_groups.png}
\end{frame}

\begin{frame}{Appendix: couple sample child composition}
\small
\centering
\safetableinput{\tabdir/sample_table_youngest_child_binned_all.tex}
\end{frame}

\begin{frame}{Appendix: couple workoutside composition over time}
\centering
\safeincludegraphics[width=0.92\linewidth,height=0.78\textheight,keepaspectratio]{\figdir/couple_workoutside_covid_wave_share.pdf}
\end{frame}

\begin{frame}{Appendix: couple WFH-some composition over time}
\centering
\safeincludegraphics[width=0.92\linewidth,height=0.78\textheight,keepaspectratio]{\figdir/couple_wfh_some_covid_wave_share.pdf}
\end{frame}

\end{document}
